\documentclass[10pt,a4paper]{amsart}
\usepackage[margin=19mm]{geometry}
\usepackage{amsmath,amssymb,mathtools}
\usepackage[scr=rsfs,cal=euler]{mathalfa}
\usepackage{xfrac}
\usepackage[unicode,hidelinks,bookmarks=false]{hyperref}
\newcommand{\ie}{i.e.\ }
\newcommand{\cf}{cf.\ }
\newcommand{\resp}{respectively\ }
\newcommand{\Subsection}{\subsection}
\providecommand{\nobreakdash}{\nobreak\hbox{-}}
\setlength{\emergencystretch}{2em}
\multlinegap0pt
\usepackage[shortlabels]{enumitem}
\usepackage{amssymb}
\usepackage{xcolor}
\usepackage{tikz}
\newtheorem{thm}{Theorem}
\newtheorem{cor}{Corollary}
\newcommand{\CFI}[1]{\CFIsym(#1)}
\newcommand{\CFIsym}{\mathsf{CFI}}
\newcommand{\FO}{\ensuremath{\mathrm{FO}}}
\newcommand{\GraphClass}{\mathcal{K}}
\newcommand{\IFPC}{\ensuremath{\mathrm{IFPC}}}
\newcommand{\IFPCWSC}{\ensuremath{\mathrm{IFPC{+}WSC}}}
\newcommand{\IFPCWSCI}{\ensuremath{\mathrm{IFPC{+}WSC{+}I}}}
\newcommand{\Iof}[1]{\mathrm{I}(#1)}
\newcommand{\PTime}{\textsc{Ptime}}
\newcommand{\Struct}{\mathfrak{A}}
\newcommand{\StructA}{\mathfrak{A}}
\newcommand{\StructB}{\mathfrak{B}}
\newcommand{\WSCIof}[1]{\mathrm{WSCI}(#1)}
\newcommand{\WSCof}[1]{\mathrm{WSC}(#1)}
\newcommand{\formA}{\Phi}
\newcommand{\interpret}{\Theta}
\newcommand{\iop}[2]{\mathsf{I}(#1;#2)}
\newcommand{\iso}{\cong}
\newcommand{\kequiv}[1]{\simeq^#1_{\mathcal{C}}}
\newcommand{\nat}{\mathbb{N}}
\newcommand{\reduct}[2]{#1 \upharpoonright #2}
\newcommand{\set}[1]{\lbrace #1 \rbrace}
\newcommand{\spleq}{\preceq}
\title{Witnessed Symmetric Choice and Interpretations in Fixed-Point Logic with Counting}
\author{Moritz Lichter}
\date{Source: arXiv:2210.07869v8 (CC BY 4.0)}
\begin{document}
\maketitle

\noindent\emph{Source and licence.} Witnessed Symmetric Choice and Interpretations in Fixed-Point Logic with Counting. Version arXiv:2210.07869v8 (CC BY 4.0), distributed by arXiv under the Creative Commons Attribution 4.0 International licence, http://creativecommons.org/licenses/by/4.0/, as recorded in the arXiv licence metadata for this exact version and shown on the abstract page https://arxiv.org/abs/2210.07869v8. Full licence notice: \texttt{licenses/arxiv-2210.07869-CC-BY-4.0.txt}.\par\smallskip

\noindent\textit{Editorial scope.} Counted unit: the article's thm thm:wsc-le-wsci (source0.tex lines 426--429). The article's complete original proof of it is delivered with it (source0.tex lines 4715--4746, one \texttt{proof} environment of 32 lines, which the article places away from the statement). No mathematics is written, repaired or re-derived: the statement and the proof are the article's own lines, in the article's own notation, and no retained line is substituted or re-flowed.  Retained uncounted as the source-local context the proof needs: lines 439--443; lines 2580--2591; lines 3580--3583; lines 4641--4655.  Nothing was omitted from any retained span, so the package carries no omission record.  No retained line contains a citation, so the wrapper carries no bibliography and no bibliography entry.  No retained line contains a figure environment, an image-inclusion command or drawing code, so no image is delivered and the package declares no asset dependency.  Editorial formatting only, in addition: an AMS \texttt{amsart} preamble that restates the article's own declarations and macros used by the retained lines, namely the environments \texttt{thm}, \texttt{cor} on separate counters, together with the standard packages those lines need.  The retained environments print in this document as Theorem 1, Corollary 1 in order of appearance, whereas the article prints the same items with the numbering of its own full document.  The complete native source of the article as distributed in the arXiv e-print for this exact version is bundled unchanged as \texttt{sources/132063/source0.tex}.
\begin{thm}
	\label{thm:canonize-CFI if-base}
	If \IFPCWSCI{} canonizes a class of colored base graphs~$\GraphClass$ (closed under individualization),
	then \IFPCWSCI{} canonizes the class of CFI graphs $\CFI{\GraphClass}$ over~$\GraphClass$.
\end{thm}

\begin{cor}
	\label{cor:canonize-CFI if-base-wsci}
	Let $\GraphClass$ be a class of base graphs.
	\begin{enumerate}
		\item If $L$ distinguishes $2$-orbits of $\GraphClass$,
		then $\CFI{\GraphClass}$ is ready for individualization in 
		$\Iof{L}$.
		\item If $\CFI{\GraphClass}$ is ready for individualization in $L$, then $\WSCof{L}$ defines a canonization for $\CFI{\GraphClass}$.
		\item If $L$ distinguishes $2$-orbits of $\GraphClass$,
		then $\WSCIof{L}$ defines a canonization for $\CFI{\GraphClass}$.
	\end{enumerate}
\end{cor}

\begin{cor}
	\label{cor:ifpc-le-wsc}
	$\IFPC{} < \IFPCWSC$.
\end{cor}

\begin{thm}
	\label{thm:asymmetric-structures-with-CFI graphs}
	There is an \FO{}-interpretation $\interpret$
	and, for every $k \in \nat$, a pair of ternary $\set{R,\spleq}$-structures $(\StructA_k, \StructB_k)$
	such that
	\begin{enumerate}[label=(\alph*)]
		\item $\spleq$ is a total preorder on $\StructA$ and $\StructB$,
		\item $\StructA_k$, $\StructB_k$, $\reduct{\StructA_k}{R}$, and $\reduct{\StructB_k}{R}$ are asymmetric,
		\item $\StructA_k \kequiv{k} \StructB_k$,
		\item $\StructA_k \not\iso \StructB_k$, and
		\item $\interpret(\StructA_k)$ and $\interpret(\StructB_k)$
		are non-isomorphic CFI graphs over the same ordered base graph.
	\end{enumerate}
	The interpretation~$\interpret$ is one-dimensional and equivalence-free.
\end{thm}

\begin{thm}
	\label{thm:wsc-le-wsci}
	$\IFPC{} < \IFPCWSC{} < \IFPCWSCI{} \leq \PTime$.
\end{thm}

\begin{proof}[Proof of Theorem~\ref{thm:wsc-le-wsci}]
	The first assertion $\IFPC{} < \IFPCWSC{}$ is Corollary~\ref{cor:ifpc-le-wsc}.
	We now show $\IFPCWSC < \IFPCWSCI{}$.
	Consider the class of $\set{R,\spleq}$-structures $\GraphClass$ given by Theorem~\ref{thm:asymmetric-structures-with-CFI graphs}.
	
	We argue that $\IFPCWSC{}=\IFPC$ on $\GraphClass'$.
	Because the structures are asymmetric, there are only singleton orbits.
	Hence, choice formulas have to define singleton choice-sets. Otherwise, it is not possible to witness them as orbits.
	Hence, choosing becomes useless and can be simulated by (non-WSC)-fixed-point operators: If the choice-set is indeed a singleton,
	its unique member is definable.
	Otherwise evaluate to false (because the WSC-fixed point operator evaluates to false if the choices were not witnessed successfully).


	Let $\interpret$ be the \FO-interpretation
	extracting the CFI graphs from $\GraphClass$-structures provided by Theorem~\ref{thm:asymmetric-structures-with-CFI graphs}.
	For every $k\in \nat$, there are two non-isomorphic structures
	$\StructA_k \kequiv{k} \StructB_k$ in~$\GraphClass$
	such that $\interpret$ evaluates on $\StructA_k$ and $\StructB_k$ to an even and an odd CFI graph, respectively.
	Consequently, we call $\Struct_k$ even and $\StructB_k$ odd.
	Since $\StructA_k \kequiv{k} \StructB_k$ for all $k\in\nat$,
	\IFPC{} does not distinguish even from odd $\GraphClass$\nobreakdash-structures.
	We show  that \IFPCWSCI{} distinguishes odd and even $\GraphClass$-structures.
	The CFI-query on ordered base graphs is definable in \IFPCWSCI{} by Theorem~\ref{thm:canonize-CFI if-base}.
	By Corollary~\ref{cor:canonize-CFI if-base-wsci},
	there is actually a $\WSCIof{\IFPC{}}=\WSCof{\IFPC}$-formula~$\formA_{\text{CFI}}$
	defining the CFI query on ordered base graphs.
	Then the $\Iof{\WSCof{\IFPC}}$-formula
	$\iop{\interpret}{\formA_{\text{CFI}}}$
	distinguishes $\StructA_k$ and $\StructB_k$ for all $k$: the interpretation $\interpret$ reduces by Theorem~\ref{thm:asymmetric-structures-with-CFI graphs}
	the problem of distinguishing even and odd $\GraphClass$~\nobreakdash structures to the  CFI query over ordered base graphs,
	which is defined by~$\formA_{\text{CFI}}$.
\end{proof}

\end{document}
